% Generated draft: requires mathematical and source review.
\documentclass[11pt]{article}
\usepackage[margin=25mm]{geometry}
\usepackage{fontspec}
\setmainfont{Latin Modern Roman}
\usepackage{amsmath,amssymb,mathtools,unicode-math}
\setmathfont{Latin Modern Math}
\usepackage{enumitem,needspace,xurl,hyperref}
\setlength{\parindent}{0pt}
\setlength{\parskip}{5pt}
\setlength{\emergencystretch}{3em}
\newcommand{\N}{\mathbb{N}}
\newcommand{\Z}{\mathbb{Z}}
\newcommand{\Q}{\mathbb{Q}}
\newcommand{\R}{\mathbb{R}}
\newcommand{\C}{\mathbb{C}}
\newcommand{\F}{\mathbb{F}}
\newcommand{\A}{\mathbb{A}}
\newcommand{\PP}{\mathbb{P}}
\newcommand{\E}{\mathbb{E}}
\newcommand{\eps}{\varepsilon}
\newcommand{\dd}{\,\mathrm{d}}
\begin{document}
\section*{UnsolvedMath ID: 20002296}
\textbf{Source problem code:} catalogue entry
\subsection{Definitions and mathematical statement}
Videla asks: ``Is there an existential analogue of Robinson's \(Q\)?'' Robinson arithmetic \(Q\) is a finitely axiomatized essentially undecidable arithmetic theory. Essential undecidability means that every consistent recursively axiomatizable extension is undecidable. The source does not say whether ``existential'' modifies the axioms, the fragment whose decision problem is considered, or an interpretation transferring undecidability. There is therefore no unique recovered theorem to prove or refute.

The supplied aid gives two precise no-go readings. First, in any fixed finite finitary first-order signature, every satisfiable finite set \(T\) of purely existential sentences has a finite model. Here each sentence has a finite existential prefix followed by a quantifier-free formula, with no tacit universal background axioms. A finite set of witnesses and all evaluated subterms gives a finite partial structure, whose operations can be extended arbitrarily while preserving the finitely many atomic and negated-atomic constraints. The resulting finite model has decidable complete theory. Consequently such \(T\) cannot be an essentially undecidable Robinson-like base or force every model to have an undecidable existential theory. This obstruction need not apply relative to a separately imposed theory excluding those finite models.

Second, use the pure ring language \(\{0,1,+,-,\cdot\}\) and all fields as the ambient class. If any set \(T\) of existential ring sentences, even infinite, has a field model \(K\), then its algebraic closure also satisfies \(T\): field embeddings preserve quantifier-free polynomial equalities and inequalities and hence existential witnesses. The algebraic closure has the decidable complete theory \(\mathrm{ACF}_p\), where \(p\) is its characteristic, including \(p=0\). Thus these existential axioms cannot force undecidability across all field models. An expansion by further named structure or a restricted class not closed under algebraic extension is a different problem.

The aid also reports a finite-field refinement: every single existential sentence in the pure ring language with a field model has a finite-field model. For a satisfied quantifier-free branch, keep finitely many witnesses and invert the finitely many polynomial values required nonzero; a finite residue field of the resulting nonzero finite-type \(\Z\)-algebra preserves those conditions. A finite set of existential sentences can be combined into one; an infinite set need not have a common finite-field model.

These obstructions address explicit readings about existential axiom bases. They do not decide an analogue defined through positive-existential interpretations, undecidability of a designated existential fragment, or a narrower geometric/arithmetic class. The aid found no authoritative reconstruction choosing among those possibilities, and does not settle the source question as a whole.
\subsection{Short English statement}
The phrase “existential analogue of Robinson's \(Q\)” is not specified enough to determine one problem. The supplied aid rules out finite existential axiom bases in unrestricted finite signatures, and existential ring axiom bases relative to all fields. Other interpretation-based or restricted-class readings remain distinct.
\subsection{Sources}
\begin{itemize}
\item Videla, Question 12, page 3, AIM workshop \emph{Definability and Decidability Problems in Number Theory}, 9--13 September 2013, moderated by Thomas Scanlon, notes by Jeroen Demeyer: \url{https://aimath.org/pastworkshops/definabilityinntproblems.pdf}.
\item Local registry \texttt{data/unsolvedmath/problems.json}, permanent ID 20002296, code \texttt{catalogue entry}; canonical \texttt{aim-logic-notes.json}, \texttt{notes[71]}; updated 14 August 2026. Clean statement null; statement status reconstructed unverified. The original literature field is empty.
\item No detailed external reference or authoritative intended definition is supplied. Attempt 1: invalid statement, problem status at run unknown, rigor label proved, novelty confidence low; current open/open-triage labels. The machine-generated aid is not peer reviewed and requires independent verification before citation.
\end{itemize}
\end{document}
